\section{第一问：预热阶段的径向分布}
\label{sec:q1}

\subsection{常热物性与非线性水分扩散}
第一问在前$1800\unit{s}$内使用附录2参数\cite{problem}。将共用方程
\eqref{eq:heat}--\eqref{eq:boundary}中的热容量、导热率分别取为
$820\times2600\unit{J/(m^3\cdot K)}$和$0.36\unit{W/(m\cdot K)}$，
水分扩散系数取$D(C)=7\times10^{-9}\exp(-0.89/C)\unit{m^2/s}$。
本问的$D$不依赖温度，热方程也不依赖含水率，因此两个场可分别求解，
不将其称为双向热质耦合模型。水分方程仍为非线性方程；对光滑解，
散度项除$D$乘拉普拉斯算子外还含$D'(C)C_r^2$，不能删去此项。

采用以轴心和表面为真实节点的圆环有限体积离散，每个内部界面的通量在
相邻控制体中符号相反；圆柱轴心用积分极限处理，不直接计算$1/r$。
时间推进使用隐式BDF方法\cite{scipy}，在每个环境输入折点重新启动积分并
传递末态。正式径向区间数为10240，相对容差$2\times10^{-12}$，最大步长
$3\unit{s}$。规定的$1\unit{s}$是输出间隔，不是积分固定步长。
均匀初态与初始表面Robin条件在角点不相容，计算保留题给初值，允许在
$t>0$形成早期边界层，不人为降低初始表面含水率。

\subsection{规定结果与预热规律}
表\ref{tab:q1-temperature}和表\ref{tab:q1-moisture}给出规定的七个时刻、
五个径向位置。\texttt{result1.xlsx}从第1秒至第1800秒，每隔
$0.1\unit{cm}$输出21个位置的温度和干基含水率，共75600个数值；
计算不提前舍入，输出统一显示四位小数。

\begin{table}[htbp]
\centering\small
\caption{30分钟内药材中截面温度（$^\circ$C）}
\label{tab:q1-temperature}
\input{tables/q1_temperature.tex}
\end{table}

\begin{table}[htbp]
\centering\small
\caption{30分钟内药材中截面干基含水率（kg/kg）}
\label{tab:q1-moisture}
\input{tables/q1_moisture.tex}
\end{table}

$1800\unit{s}$时轴心、表面温度分别为33.5753\celsius{}和
36.7856\celsius，均低于该时刻环境温度41.5130\celsius。表面先响应，
轴心因传导而滞后。由未舍入值计算，径向温差为$3.2102\unit{K}$。
同期表面含水率降至$1.5102\unit{kg/kg}$，轴心实际含水率约为
$2.5499924\unit{kg/kg}$，四位显示仍为2.5500；因此不能把表中轴心数字
未变解释为水分数学上严格不变。

\begin{figure}[htbp]
\centering
\includegraphics[width=0.94\textwidth]{figures/q1_profiles.pdf}
\caption{第一问不同时刻的径向温度与干基含水率分布。}
\label{fig:q1-profiles}
\end{figure}

图\ref{fig:q1-profiles}显示，升温已遍及所考察的径向位置，而水分变化
主要集中在表层。末时$r=1.5\unit{cm}$处的含水率为2.3755，较初值
下降约6.84\%，$r=1.0\unit{cm}$处仍为2.5383。这里不存在严格截断的
扩散前沿；描述受影响层厚度时需要先规定含水率变化阈值。

以共同尺度$R$比较扩散时间：初始热扩散率
$\alpha=k/(\rho c_p)=1.6886\times10^{-7}\unit{m^2/s}$，
$D(C_0)=4.9377\times10^{-9}\unit{m^2/s}$，故
$R^2/\alpha=2368.89\unit{s}$、$R^2/D(C_0)=81010.11\unit{s}$，
两者相差34.20倍。这解释了预热阶段热响应快于水分响应；因失水使$D$
继续减小，初始扩散时间不是最终烘干时长。
圆柱平均须按$r\,\mathrm dr$加权，得到末时中截面平均温度
35.1679\celsius、平均含水率$2.2936\unit{kg/kg}$；它们不替代最湿点，
也不能直接当作包含端部效应的整根平均量。

\subsection{数值可信范围}
温度采用贝塞尔特征函数展开作独立参考，水分采用
$x=(r/R)^2$变换的Chebyshev配点与非线性表面条件消元作独立空间参考。
正式解与参考在全部规定输出点的最大温差为
$7.29\times10^{-9}\unit{K}$，含水率差为
$1.52\times10^{-6}\unit{kg/kg}$；全部75600个正式输出四位一致。
固定正式网格收紧时间控制后，温湿差分别约
$5.03\times10^{-10}\unit{K}$和$8.15\times10^{-11}\unit{kg/kg}$。
接近舍入分界的数值必须逐项核对，不能仅凭最大误差小于半个末位单位
推断整张表四位一致。

同物理条件的$80\times128$有限圆柱二维配对，在前$1800\unit{s}$
已检查采样点的中截面最大温湿差约为
$5.65\times10^{-8}\unit{K}$与$4.09\times10^{-9}\unit{kg/kg}$，
支持本问中截面径向近似；端面局部差仍可较大，不能据此推广到整个圆柱。
以上验证支持所选有效方程的数值结果，不消除等效湿度边界、参考密度
及无显式潜热假设引入的物理不确定性。
\FloatBarrier
